\section{Tier 1 Visual Adjudication}

\subsection{Purpose}

With \texttt{692} now located in the local Fuls 2013 positional-analysis article, the project has both visual sides for all four Tier 1 distributional candidates. This section converts evidence collection into a first visual adjudication. The goal is not to merge signs. The goal is to separate plausible allograph candidates from signs that merely share a positional or functional profile.

The adjudication follows the visual-caution protocol: catalog shape, positional behavior, artifact context, direction, and object-level independence must all support a merge before the inventory is changed \citep{daggumati2021allographs}. Fuls' positional analysis is especially useful here because it argues that visually similar signs can still be separate graphemes when their positional behavior differs \citep{fuls2013positional}.

\subsection{Provisional Decisions}

\begin{table}[htbp]
\centering
\begin{tabular}{p{0.13\textwidth}p{0.34\textwidth}p{0.38\textwidth}}
\toprule
\textbf{Pair} & \textbf{Visual assessment} & \textbf{Provisional decision} \\
\midrule
\texttt{817}/\texttt{861} & Moderate family resemblance: both are enclosed signs with angular internal marking, but the outline and inner geometry differ. & Keep distinct. Treat as a functional-cluster candidate, not a merge candidate. \\
\texttt{817}/\texttt{820} & Moderate family resemblance: both are enclosed signs, but \texttt{817} has a simpler chevron-like interior while \texttt{820} has a divided/crossed interior. & Keep distinct. Treat as a functional-cluster candidate, not a merge candidate. \\
\texttt{705}/\texttt{706} & Strong visual resemblance: same open-U/long-stem family, with \texttt{706} adding or emphasizing an inner nested stroke. & Variant-link candidate. Highest priority for artifact-level review; no merge yet. \\
\texttt{692}/\texttt{920} & Weak visual resemblance: \texttt{692} is crossed/X-shaped, while \texttt{920} is a curved single-stroke form. & Reject visual-allograph hypothesis at catalog level. Keep as a functional-cluster/control case. \\
\bottomrule
\end{tabular}
\caption{First visual adjudication of Tier 1 distributional candidates.}
\label{tab:tier1-visual-adjudication}
\end{table}

\subsection{Interpretation}

The important result is that distributional similarity is not the same as visual equivalence. The strongest allograph-style candidate is \texttt{705}/\texttt{706}. The pairs \texttt{817}/\texttt{861} and \texttt{817}/\texttt{820} look more like members of a shared initial-sign family than variants of one sign. The pair \texttt{692}/\texttt{920} is now valuable as a control case: it is distributionally strong but visually weak.

\subsection{Outputs}

\begin{itemize}
  \item \texttt{outputs/tier1\_visual\_adjudication.csv}
  \item \texttt{outputs/tier1\_visual\_adjudication.tex}
  \item \texttt{outputs/tier1\_artifact\_review\_queue\_705\_706.csv}
  \item \texttt{outputs/tier1\_artifact\_review\_summary\_705\_706.csv}
  \item \texttt{outputs/tier1\_artifact\_review\_705\_706.tex}
\end{itemize}
